\section{Marker separation on the LCG and LCGS parents}
\label{sec:marker-separation}

This section supplies a table-level relative proof for two new compositions,
\texttt{PA194LCGMS} and \texttt{PA194LCGSMS}.  Both have 194 working controls.
The marker-separation idea is the existing \texttt{PA194B} construction in
\texttt{PA194B\_PA194G\_research\_2026-11-14}; the new contribution here is its
composition with the pinned \texttt{PA194LCG} and \texttt{PA194LCGS} parents,
including their loop reassociation and guarded return shortcuts.  This is a
new allocation of recurring physical operations, with a relative proof that
is independent of the numerical-calculus and register-simulation proofs.
It is not a claim that the marker-separation mechanism itself was invented
in this investigation.

\subsection{Exact transformation}
Let $P$ be either pinned native parent.  Its initial and terminal controls
are respectively 74 and 254.  Define
\[
 f(51)=400,\qquad f(301)=88,\qquad f(q)=q
 \quad(q\notin\{51,301\}).
\]
Control 400 does not occur in either parent.  For every old row other than a
row of control 301, initially copy the row with every control name replaced
by $f$.  Then replace the row at $(88,0)$ and introduce the two rows at 51:
\[
 \begin{array}{c|cc}
       &0&1\\\hline
 88&0\mathsf L51&1\mathsf R103\\
 51&0\mathsf R122&1\mathsf R51\\
 400&0\mathsf R16&1\mathsf R400.
 \end{array}
\]
The last row is the renamed original marker scanner.  Thus all old entries
into the marker scanner now enter 400; control 51 is reserved for decoder
selection.  Exactly one old control, 301, is removed and exactly one fresh
control, 400, is introduced.  The resulting machine $N$ has exactly 194
working controls and 388 transition rows.  Breadth-first enumeration from
74 visits every working control; the delivered dense conjugacy sends 74 to
0 and 254 to the separate halt label 194.

The following parent rows are common to both inputs:
\[
\begin{array}{c|cc}
 &0&1\\\hline
301&0\mathsf R45&1\mathsf R103\\
88&0\mathsf R122&1\mathsf R103\\
122&0\mathsf L137&0\mathsf L1\\
137&0\mathsf R30&1\mathsf R70\\
30&0\mathsf R45&1\mathsf R46.
\end{array}
\]
The only incoming parent edge to 301 is the zero row of 87, which writes
zero and moves right.  The only incoming edge to 88 is the one row of 77,
which writes one and moves right.  Since the initial control is neither,
every reachable occurrence of 301 has a zero immediately left of the head;
every reachable occurrence of 88 has a one there.  This is an exact
transition-graph invariant, valid for arbitrary complete initial tapes.
It requires no assertion about which source instructions can occur.

\subsection{Four exhaustive local cases}
Write the current head coordinate as 0 and the preceding coordinate as
$-1$.  The head bit is $b$ and the preceding bit is $h$.  The complete tape
outside these two cells is arbitrary and unchanged.  The local substitution
has the following behavior:
\[
\begin{array}{ccccc}
\text{old control}&h&b&\text{new duration}&\text{common exit}\\\hline
301&0&0&5&(45,1)\\
301&0&1&1&(103,1)\\
88&1&0&3&(122,1)\\
88&1&1&1&(103,1).
\end{array}
\]
Both cells retain their original bits in all four cases.  The two nontrivial
control/head traces are
\[
\begin{split}
 (88,0)&\longrightarrow(51,-1)\longrightarrow(122,0)
 \longrightarrow(137,-1)\longrightarrow(30,0)
 \longrightarrow(45,1) \qquad(h,b)=(0,0),\\
 (88,0)&\longrightarrow(51,-1)\longrightarrow(51,0)
 \longrightarrow(122,1) \qquad(h,b)=(1,0).
\end{split}
\]
Each displayed transition reads only one of the two known cells.  No value
of any other cell can affect these computations.  When $b=1$, the common
row $1\mathsf R103$ gives the exit in one step.  Every remaining old row is
simulated in one step after applying $f$, including the marker scanner now
named 400.  The executable checker verifies all unaffected rows and all
four complete two-cell traces against the final simultaneously transformed
table.  It also enumerates all incoming edges establishing the invariant.

\subsection{Explicit clock bounds and intuitionistic transfer}
Represent a finite-support two-way tape by the usual pair of natural-number
binary words together with a head position, and totalize both machines by
making the halt configuration absorbing.  The machine step and its bounded
iteration are primitive recursive.  Apply $f$ to the control component of
a configuration and leave its tape and head unchanged; call the resulting
map $F$.

For the old trajectory $C_P(t)$, define a primitive-recursive duration
$\tau(t)\in\{1,3,5\}$: it is 5 at an old $(301,0)$ step, 3 at an old
$(88,0)$ step, and 1 otherwise.  Put
\[
 \theta(0)=0,\qquad \theta(t+1)=\theta(t)+\tau(t).
\]
The incoming-history invariant and the four local traces prove by induction
\[
 C_N(\theta(t))=F(C_P(t)),\qquad t\leq\theta(t)\leq5t.
\]
At a halt, use the totalized one-step self-loop and $\tau=1$.
The initial equality holds because $f(74)=74$ and both initial tapes are
the same.  The induction predicates are equalities and inequalities between
primitive-recursive functions.  The finite table case distinction and the
incoming-history induction are equally elementary.  Consequently this
argument is formalizable by quantifier-free induction in intuitionistic
primitive recursive arithmetic, without excluded middle for arbitrary
formulas, an unbounded search, a soundness assumption about PA, or an
assumption of PA consistency.

Let $H_P(t)$ and $H_N(t)$ mean that the corresponding totalized blank run
has reached halt by time $t$.  Since $F$ maps no working control to halt,
monotonicity of totalized halting gives the sharper, asymmetric estimates
\[
 \mathsf{iPRA}\vdash H_P(t)\to H_N(5t),\qquad
 \mathsf{iPRA}\vdash H_N(t)\to H_P(t).
\]
For the second implication, $H_N(t)$ and $t\leq\theta(t)$ give
$H_N(\theta(t))$, and the displayed configuration equality then gives
$H_P(t)$.  No reverse unbounded simulation is being assumed.
Writing $\mathrm{NH}(M):=\forall t\,\neg H_M(t)$, intuitionistic implication
reasoning now yields
\[
 \mathsf{iPRA}\vdash\mathrm{NH}(P)\leftrightarrow\mathrm{NH}(N).
\]
Therefore any established ordinary or formalized iPRA consistency
calibration for either parent transfers to its child.  This section proves
the new relative obligation.  The absolute equivalence with
$\mathrm{Con}(\mathsf{PA})$ still invokes the supplied written calibration
of the LCG/LCGS parents and their AX11H dependency.  No end-to-end Lean
kernel theorem for the new bytes is claimed here.

\subsection{Complete blank initialization}
The checker replays every step of each parent's pinned blank initialization,
matching it after between one and five child steps.  At the endpoint it
checks the entire finite set of one-cells, the head, the control, and the
blank exterior against the canonical source frame.  The logical source
entry is site 25; the head is at $-2$, the gap parameter is 48, and the
physical register file is
\[
 (0,0,0,0,0,0,2,0,0)\mathbin{+\!+}(\underbrace{1,\ldots,1}_{26}).
\]
\[
\begin{array}{lrr}
\text{child}&\text{parent steps}&\text{child steps}\\\hline
\texttt{PA194LCGMS}&57\,870&57\,970\\
\texttt{PA194LCGSMS}&57\,712&57\,812.
\end{array}
\]
These finite initialization replays are corroborating identities.  The
all-time equivalence follows from the local simulation and its induction,
not from extrapolation of either finite run.
